% ============================================================
% droop_first_tests.tex
% Rapport: Foerste tester av den ytre frekvensstoette-droopen (Fase 1)
% Frittstaaende dokument for oversikt (student + veileder).
% Bygg:  pdflatex droop_first_tests.tex   (kjoer fra denne mappa)
% Figurer hentes fra diagnostic_plots via \graphicspath nedenfor.
%
% NB (avklart 2026-07-28):
%   - "Ti" = UIC-ens interne-spennings frekvensfilter T_filter (UIC.py):
%     Told = 0.1 s, Tnew = 0.01 s. MERK: T_filter paavirker kun dynamikken naar
%     perfect_tracking = 1; naavaerende build_model har perfect_tracking = 0, saa
%     disse figurene maa vaere kjoert med perfect_tracking = 1 (bekreft).
%   - K_droop = 0.50 pu/Hz (drivers default; build_model har 0.75 men overstyres).
%     headroom = 0.05 pu (= 1 MW paa 20 MVA UIC-base), f_nom = 50 Hz.
% ============================================================
\documentclass[11pt]{article}

\usepackage[a4paper,margin=25mm]{geometry}
\usepackage{amsmath}
\usepackage{amssymb}
\usepackage{graphicx}
\usepackage{booktabs}
\usepackage{siunitx}
\usepackage{hyperref}

\DeclareSIUnit{\pu}{pu}

% Figurer ligger utenfor docs/, i resultat-treet. Juster ved behov.
\graphicspath{{../../casestudies/dyn_sim/results/diagnostic_plots/}}

\title{First Tests of the Outer Wind-Turbine\\ Frequency-Support Droop on the LEOGO Grid}
\author{Summer project 2026 --- work log for supervisor review}
\date{\today}

\begin{document}
\maketitle

\begin{abstract}
\noindent This note documents the \emph{first} set of tests carried out with the
outer wind-turbine (WT) frequency-support droop, run on the islanded LEOGO
grid. The purpose is a shared overview of \emph{what was actually done} and
\emph{what the results show}, not a full tuning study. A de-loaded reduced WT,
coupled to the grid through the Unified Impedance Converter (UIC), is compared
with and without the droop for a series of grid load steps. The tests confirm
that the reserved headroom is released as extra active power during
under-frequency, and that a controller retuning (labelled \texttt{OldTi} vs
\texttt{NewTi}) mainly improves the transient damping without changing the
steady frequency-support gain. This establishes the Phase-1 baseline for the
later drivetrain- and tower-loading studies.
\end{abstract}

% ------------------------------------------------------------
\section{Purpose and scope}
\label{sec:dft-purpose}
These were the first experiments in the project. The question was simple: can a
de-loaded wind turbine, connected through the grid-forming UIC, release its
reserve to support the frequency of the small islanded LEOGO network when the
grid load suddenly increases? The tests were designed to (i) verify the intended
droop action, (ii) record the full electromechanical response (converter,
drivetrain, pitch), and (iii) provide a clean baseline against which the later
mechanical-loading trade-off is measured. This is a work log, so the emphasis is
on the test matrix and the qualitative behaviour rather than on parameter
optimisation.

% ------------------------------------------------------------
\section{Model and experimental setup}
\label{sec:dft-setup}
The plant is a reduced wind-turbine model (two-mass drivetrain, pitch
controller, and maximum-power-tracking generator-torque law) coupled to the
islanded LEOGO network through the UIC grid-forming converter. The grid
frequency seen by the turbine is the centre-of-inertia (COI) frequency of the
synchronous machines, i.e.\ an inertia-weighted average of their speeds, fed
back to the WT droop as \(f_{grid}\).

Each test runs a \emph{matched pair}:
\begin{itemize}
  \item \textbf{droop off} --- de-loaded, no frequency response;
  \item \textbf{droop on} --- same de-loading, frequency response enabled.
\end{itemize}
Both cases are de-loaded by the \emph{same} headroom reserve, so the pre-event
operating point is identical and the only difference is the droop action. The
control law (documented in \texttt{outer\_droop.tex}) is
\begin{align}
  P_{base} &= \operatorname{clip}\!\big(P_{avail}-P_{hr},\,0,\,P_{avail}\big),\\
  \Delta P &= K_{droop}\,(f_{nom}-f_{grid}),
\end{align}
so under-frequency (\(f_{grid}<f_{nom}\)) commands additional active power
\(\Delta P>0\) up to the reserved headroom \(P_{hr}\). The driver script is
\texttt{test\_WT\_LEOGO\_droop\_comparison\_diagnostics.py}.

The droop parameters, referred to the \SI{20}{\mega\voltampere} UIC base, are
\(f_{nom}=\SI{50}{\hertz}\), a headroom reserve
\(P_{hr}=\SI{0.05}{\pu}\;(\approx\SI{1}{\mega\watt})\), and a droop gain
\(K_{droop}=\SI{0.50}{\pu\per\hertz}\) (the comparison-driver default;
\texttt{build\_model} in \texttt{test\_WT\_LEOGO\_sim.py} configures
\SI{0.75}{\pu\per\hertz}, but the driver overrides it). As shown below, for the
\SI{5}{\mega\watt} step the droop saturates at the headroom, so the steady
response is set by \(P_{hr}\) and is insensitive to this gain.

% ------------------------------------------------------------
\section{Test matrix}
\label{sec:dft-matrix}
A grid load step is applied at \(t\approx\SI{5}{\second}\), held for about
\SI{50}{\second}, in a run of roughly \SI{200}{\second}. For every case four time
windows are exported: the full timeline, a zoom around load application, a zoom
around load removal, and one additional window. Two step magnitudes, two step
profiles, and two controller tunings are combined (Table~\ref{tab:dft-matrix}).
A separate set of runs uses a sinusoidal load (\SI{1.2}{\hertz}, small
amplitude) as a forced-response probe.

\begin{table}[htbp]
  \centering
  \caption{Load-step test cases. Each case is run droop-off and droop-on, and
           stored under \texttt{diagnostic\_plots/}.}
  \label{tab:dft-matrix}
  \begin{tabular}{llll}
    \toprule
    Step [\si{\mega\watt}] & Profile & Tuning & Folder \\
    \midrule
    1 & abrupt & OldTi & \texttt{1MWAbrupt\_OldTi}\\
    1 & abrupt & NewTi & \texttt{1MWAbrupt\_NewTi}\\
    1 & smooth & OldTi & \texttt{1MWSmooth\_OldTi}\\
    1 & smooth & NewTi & \texttt{1MWSmooth\_NewTi}\\
    5 & abrupt & OldTi & \texttt{5MWAbrupt\_OldTi}\\
    5 & abrupt & NewTi & \texttt{5MWAbrupt\_NewTi}\\
    5 & smooth & OldTi & \texttt{5MWSmooth\_OldTi}\\
    5 & smooth & NewTi & \texttt{5MWSmooth\_NewTi}\\
    \midrule
    \multicolumn{4}{l}{Sinusoidal-load probes: \texttt{Sine\_1.2Hz*}, \texttt{SInus\_load}}\\
    \bottomrule
  \end{tabular}
\end{table}

\noindent Here \emph{abrupt} is an (essentially) ideal step and \emph{smooth} is
a ramped step. \texttt{OldTi} and \texttt{NewTi} denote two settings of the UIC
internal-voltage frequency-filter time constant \(T_{filter}\)
(\texttt{OldTi}: \(\SI{0.1}{\second}\); \texttt{NewTi}: \(\SI{0.01}{\second}\));
its effect is discussed in Section~\ref{sec:dft-tuning}.

% ------------------------------------------------------------
\section{Diagnostics recorded}
\label{sec:dft-diag}
Fifteen signal families are logged per case and exported as comparison plots
(droop-off vs droop-on). Grouped, they are:
\begin{enumerate}
  \item converter active power, and the isolated droop contribution;
  \item the power chain \(P_{aero}\rightarrow P_{e}\rightarrow P_{ref}\);
  \item torque balance and net (accelerating) torque;
  \item drivetrain speeds (rotor and generator);
  \item pitch error, pitch position, pitch rate;
  \item converter terminal voltage, current magnitude and angle;
  \item bus power vs WT-torque power, and bus-power tracking.
\end{enumerate}

% ------------------------------------------------------------
\section{Results and observations}
\label{sec:dft-results}

\subsection{Frequency support is delivered}
\label{sec:dft-support}
Figure~\ref{fig:dft-power} (5\,MW abrupt, NewTi) shows the converter active
power. Without droop it stays near \SI{0.140}{\pu} and only shows brief
transients at load application and removal. With droop it rises to about
\SI{0.1497}{\pu} while the grid is under-frequency, and returns after the load
clears. Figure~\ref{fig:dft-contrib} isolates the droop contribution: the extra
active power builds up to roughly \SI{+0.0097}{\pu} (system base) during the
event and decays after clearing. This is the reserve being released from the
headroom to arrest the frequency drop --- the intended behaviour. Note that this
plateau (\(\approx\SI{0.0097}{\pu}\) on the \SI{100}{\mega\voltampere} system
base, i.e.\ \(\approx\SI{1}{\mega\watt}\)) coincides with the full headroom
reserve \(P_{hr}\): for a \SI{5}{\mega\watt} step the droop reaches its reserve
limit, so the steady contribution is set by the headroom rather than by
\(K_{droop}\) (the gain governs the ramp rate and the response to smaller steps,
e.g.\ the \SI{1}{\mega\watt} cases).

\begin{figure}[htbp]
  \centering
  \includegraphics[width=0.82\linewidth]{5MWAbrupt_NewTi/event_full/01_uic_active_power_comparison.png}
  \caption{Converter active power, droop-off vs droop-on (5\,MW abrupt step,
           NewTi). With droop the turbine releases extra power during the
           under-frequency window.}
  \label{fig:dft-power}
\end{figure}

\begin{figure}[htbp]
  \centering
  \includegraphics[width=0.82\linewidth]{5MWAbrupt_NewTi/event_full/02_droop_contribution_comparison.png}
  \caption{Isolated droop contribution (5\,MW abrupt step, NewTi): about
           \SI{+0.0097}{\pu} of additional active power during the load event,
           decaying after the load is cleared.}
  \label{fig:dft-contrib}
\end{figure}

\subsection{Matched operating point}
\label{sec:dft-matched}
Because both cases start from the identical de-loaded point, the difference
between the two curves is purely the droop action. This clean isolation is the
reason the tests use a headroom reserve even in the droop-off case.

\subsection{Effect of the controller tuning: OldTi vs NewTi}
\label{sec:dft-tuning}
Comparing the same 5\,MW abrupt step with the two tunings
(Figure~\ref{fig:dft-oldti} vs Figure~\ref{fig:dft-power}) shows that the
retuning mainly changes the \emph{damping of the converter active-power
transient} at load application and removal. With \texttt{OldTi} the power rings
for several seconds after each step, whereas with \texttt{NewTi} it settles
quickly. The two tunings differ only in the UIC internal-voltage frequency-filter
time constant \(T_{filter}\) (\(\SI{0.1}{\second}\rightarrow\SI{0.01}{\second}\)):
the shorter constant lets the converter's internal voltage track the frequency
signal faster and damps the transient. Importantly, the \emph{steady} droop
contribution (\(\approx\SI{0.0097}{\pu}\)) is essentially unchanged --- the
retuning improves the transient behaviour without altering the frequency-support
gain.

\begin{figure}[htbp]
  \centering
  \includegraphics[width=0.82\linewidth]{5MWAbrupt_OldTi/event_full/01_uic_active_power_comparison.png}
  \caption{Converter active power with the \texttt{OldTi} tuning (5\,MW abrupt
           step). The transient at load on/off is markedly less damped than with
           \texttt{NewTi} (Figure~\ref{fig:dft-power}); the steady droop support
           is the same.}
  \label{fig:dft-oldti}
\end{figure}

\subsection{Step size and profile}
\label{sec:dft-size}
A \SI{5}{\mega\watt} step produces a larger frequency deviation --- and hence a
larger droop contribution --- than a \SI{1}{\mega\watt} step. The \emph{smooth}
(ramped) profile reduces the transient spike relative to the \emph{abrupt} step
while giving the same steady response. Across all cases the bus-power tracking
plot confirms that the converter bus power follows its reference closely
throughout the event.

% ------------------------------------------------------------
\section{Status and open items}
\label{sec:dft-status}
These first tests establish the Phase-1 baseline: a de-loaded wind turbine,
connected through the UIC, provides measurable and controllable frequency
support on the islanded LEOGO grid, and the \texttt{NewTi} retuning cleans up the
converter transient without sacrificing support. Open items:
\begin{itemize}
  \item consistency check: \(T_{filter}\) only affects the dynamics when the UIC
        \texttt{perfect\_tracking} feedback is enabled; the current
        \texttt{build\_model} has it disabled, so these figures were presumably
        produced with it enabled --- worth confirming before archiving;
  \item the later phases quantify the drivetrain- and tower-loading \emph{cost}
        of providing this frequency support.
\end{itemize}

\end{document}
